% =============================================================================
% Chapitre 4 : Architecture du système
% =============================================================================
\chapter{Architecture du système}
\label{ch:architecture}

Ce chapitre présente l'architecture globale, les flux de données et les choix de conception du pipeline de modération.

\section{Vue d'ensemble}

Le système suit une architecture en couches avec séparation claire des responsabilités (figure~\ref{fig:architecture_overview}) : couche clients (React, Streamlit, consommateurs API), couche API FastAPI, couche ML (modèles texte/image et routeur de confiance), couche services (file de modération, registre de modèles, analytique, réentraînement) et couche données (SQLAlchemy, migrations Alembic, journal d'audit), adossée à PostgreSQL, aux checkpoints de modèles et au suivi MLflow.

\begin{figure}[H]
\centering
\scalebox{0.8}{\begin{tikzpicture}[
  node distance=0.6cm and 0.3cm,
  layer/.style={
    rectangle, rounded corners=4pt,
    draw=gray!60, fill=gray!5,
    minimum width=13cm, minimum height=1.2cm,
    align=center, font=\small
  },
  component/.style={
    rectangle, rounded corners=3pt,
    draw=blue!70, fill=blue!8,
    minimum width=2.4cm, minimum height=0.9cm,
    align=center, font=\footnotesize
  },
  db/.style={
    cylinder, cylinder uses custom fill,
    cylinder end fill=gray!20, cylinder body fill=gray!10,
    draw=gray!60, minimum width=1.8cm, minimum height=0.9cm,
    align=center, font=\footnotesize, shape border rotate=90
  },
  arrow/.style={-{Stealth[length=2.5mm]}, thick, draw=gray!70},
  every node/.style={font=\small}
]

% Layer 1: Client Layer
\node[layer] (client) {
  \textbf{Couche Clients}\\
  \scriptsize{Frontend React (Vite) \quad|\quad Tableau de bord Streamlit \quad|\quad Clients API}
};

% Layer 2: API Layer
\node[layer, below=of client] (api) {
  \textbf{Couche API (FastAPI)}\\
  \scriptsize{POST /moderate/text \quad|\quad POST /moderate/image \quad|\quad POST /moderate/content \quad|\quad GET /moderation\_queue \quad|\quad POST /review}
};

% Layer 3: ML Layer
\node[layer, below=of api] (ml) {
  \textbf{Couche ML}\\
  \scriptsize{TextModerationModel (BERT) \quad|\quad ImageModerationModel (CLIP + NudeNet) \quad|\quad ConfidenceRouter}
};

% Layer 4: Services Layer
\node[layer, below=of ml] (svc) {
  \textbf{Couche Services}\\
  \scriptsize{QueueService \quad|\quad ModelRegistry \quad|\quad AnalyticsService \quad|\quad RetrainPipeline}
};

% Layer 5: Data Layer
\node[layer, below=of svc] (data) {
  \textbf{Couche Données}\\
  \scriptsize{SQLAlchemy ORM \quad|\quad Migrations Alembic \quad|\quad AuditLog}
};

% Database
\node[db, below right=0.4cm and 2cm of data.south west, anchor=north] (postgres) {PostgreSQL\\16};

% Model storage
\node[component, below=0.4cm of data.south, anchor=north] (models) {Checkpoints\\de modèles};

% MLflow
\node[component, below left=0.4cm and 2cm of data.south, anchor=north] (mlflow) {Suivi\\MLflow};

% Arrows
\draw[arrow] (client) -- (api);
\draw[arrow] (api) -- (ml);
\draw[arrow] (ml) -- (svc);
\draw[arrow] (svc) -- (data);
\draw[arrow] (data) -- (postgres);
\draw[arrow] (data) -- (models);
\draw[arrow] (data) -- (mlflow);

\end{tikzpicture}}
\caption{Architecture haut niveau du système : pile à cinq couches et dépendances externes.}
\label{fig:architecture_overview}
\end{figure}

\section{Architecture de déploiement}

Le système est déployé comme une pile Docker Compose de six conteneurs indépendamment constructibles et redémarrables (figure~\ref{fig:deployment}).

\begin{figure}[H]
\centering
\scalebox{0.8}{\begin{tikzpicture}[
  node distance=0.8cm and 0.6cm,
  svc/.style={
    rectangle, rounded corners=5pt,
    draw=blue!60, fill=blue!5,
    minimum width=2.8cm, minimum height=1.6cm,
    align=center, font=\footnotesize, thick
  },
  net/.style={
    dashed, draw=gray!50, rounded corners=8pt,
    inner sep=8pt
  },
  arrow/.style={-{Stealth[length=2mm]}, thick, draw=gray!60},
  label/.style={font=\tiny, fill=white, inner sep=1pt, midway}
]

% Network boundary
\node[net, minimum width=13cm, minimum height=8cm] (netbox) {};

% PostgreSQL
\node[svc, fill=blue!15] (postgres) at (-4.5, 2.5) {PostgreSQL\\16\\{\tiny\texttt{port: 5432}}};

% MLflow
\node[svc, fill=orange!10] (mlflow) at (0, 2.5) {MLflow\\Server\\{\tiny\texttt{port: 5000}}};

% API
\node[svc, fill=green!10] (api) at (4.5, 2.5) {Backend\\FastAPI\\{\tiny\texttt{port: 8000}}};

% Frontend
\node[svc, fill=purple!10] (frontend) at (-4.5, -1) {Frontend\\React\\{\tiny\texttt{port: 3000}}};

% Dashboard
\node[svc, fill=teal!10] (dashboard) at (0, -1) {Tableau de bord\\Streamlit\\{\tiny\texttt{port: 8501}}};

% Retrain
\node[svc, fill=red!10] (retrain) at (4.5, -1) {Planificateur\\de réentraînement\\{\tiny basé cron}};

% Arrows showing dependencies
\draw[arrow] (api) -- node[label, above] {SQL} (postgres);
\draw[arrow] (api) -- node[label, above] {tracking} (mlflow);
\draw[arrow] (dashboard) -- (postgres);
\draw[arrow] (dashboard) -- (api);
\draw[arrow] (frontend) -- (api);
\draw[arrow] (retrain) -- (postgres);
\draw[arrow] (retrain) -- (mlflow);

% Network label
\node[below, font=\small, gray] at (0, -3.5) {\texttt{moderation-net (bridge)}};

\end{tikzpicture}}
\caption{Architecture de déploiement Docker Compose : six conteneurs et leurs communications.}
\label{fig:deployment}
\end{figure}

\section{Flux de données}

Pour une requête de modération, le flux est identique pour les deux modalités : entrée $\to$ prédiction (BERT pour le texte avec sorties sigmoïdes ; CLIP + NudeNet pour l'image avec score NSFW) $\to$ routage par confiance (résolution automatique si le contenu est propre ou confiance élevée ; file des modérateurs si confiance moyenne ; revue manuelle sinon) $\to$ persistance PostgreSQL $\to$ réponse. Les décisions image suivent trois seuils : SAFE si score $< 0{,}30$, NSFW si score $> 0{,}70$, REVIEW entre les deux.

\section{Schéma de base de données}

La persistance repose sur six tables normalisées (figure~\ref{fig:er_diagram}), suivant le principe d'audit append-only : la table \texttt{audit\_logs} enregistre tout événement modifiant l'état et n'est jamais mise à jour ni supprimée.

\begin{figure}[H]
\centering
\scalebox{0.85}{\begin{tikzpicture}[
  entity/.style={
    rectangle, rounded corners=4pt,
    draw=blue!60, fill=blue!5,
    minimum width=3cm, minimum height=0.6cm,
    align=left, font=\scriptsize, inner sep=4pt
  },
  pk/.style={font=\scriptsize\bfseries},
  fk/.style={font=\scriptsize, text=purple!70},
  arrow/.style={-{Stealth[length=2mm]}, thick, draw=gray!50},
  label/.style={font=\tiny, fill=white, inner sep=1pt},
  every node/.style={font=\scriptsize}
]

% User
\node[entity] (user) at (0, 4) {
  \textbf{User}\\
  \pk{id} (PK)\\
  username, email\\
  hashed\_password\\
  role (enum)\\
  is\_active
};

% ModerationRequest
\node[entity] (request) at (5, 4) {
  \textbf{ModerationRequest}\\
  \pk{id} (PK)\\
  content\_type (enum)\\
  text\_content\\
  image\_path\\
  status (enum)\\
  \fk{submitted\_by\_id (FK $\to$ users)}
};

% Prediction
\node[entity] (prediction) at (5, 1.5) {
  \textbf{Prediction}\\
  \pk{id} (PK)\\
  \fk{moderation\_request\_id (FK)}\\
  modality (enum)\\
  model\_version\\
  top\_label, top\_score\\
  label\_scores (JSON)
};

% ModeratorReview
\node[entity] (review) at (0, 1.5) {
  \textbf{ModeratorReview}\\
  \pk{id} (PK)\\
  \fk{moderation\_request\_id (FK)}\\
  \fk{moderator\_id (FK $\to$ users)}\\
  decision (enum)\\
  corrected\_labels (JSON)
};

% ModelVersion
\node[entity] (modelversion) at (-5, 1.5) {
  \textbf{ModelVersion}\\
  \pk{id} (PK)\\
  model\_type (enum)\\
  version\_tag\\
  checkpoint\_path\\
  metrics (JSON)\\
  is\_active
};

% AuditLog
\node[entity] (auditlog) at (-5, 4) {
  \textbf{AuditLog}\\
  \pk{id} (PK)\\
  \fk{moderation\_request\_id (FK)}\\
  event\_type\\
  actor\\
  details (JSON)
};

% Relationships
\draw[arrow] (request) -- node[label, left, pos=0.4] {1:N} (user.east|-request.west);
\draw[arrow] (prediction) -- node[label, right, pos=0.3] {N:1} (request);
\draw[arrow] (review) -- node[label, left, pos=0.3] {1:1} (request.west|-review.east);
\draw[arrow] (review) -- node[label, left, pos=0.4] {N:1} (user.west|-review.north west);
\draw[arrow] (auditlog) -- node[label, left, pos=0.4] {N:1} (request.west|-auditlog.south);

\end{tikzpicture}}
\caption{Diagramme entité-association : les six tables cœur de la base et leurs relations.}
\label{fig:er_diagram}
\end{figure}

\section{Conception de l'API}

L'API REST suit les conventions RESTful avec des endpoints versionnés sous \texttt{/api/v1}.

\begin{table}[H]
\centering
\caption{Endpoints de l'API REST.}
\label{tab:api_endpoints}
\begin{tabularx}{\textwidth}{lllX}
\toprule
\textbf{Méthode} & \textbf{Endpoint} & \textbf{Tag} & \textbf{Description} \\
\midrule
POST & \texttt{/moderate/text} & Modération & Classer un texte (3 labels) \\
POST & \texttt{/moderate/image} & Modération & Classer une image (5 labels) \\
POST & \texttt{/moderate/content} & Modération & Modération combinée texte + image \\
GET & \texttt{/moderation\_queue} & RH & Lister les éléments en attente de revue \\
POST & \texttt{/review} & RH & Soumettre une décision de modérateur \\
GET & \texttt{/metrics} & Système & Statistiques opérationnelles agrégées \\
GET & \texttt{/health} & Système & Sonde de vivacité/disponibilité \\
\bottomrule
\end{tabularx}
\end{table}

\section{Configuration et patterns de conception}

Toute la configuration d'exécution est centralisée dans une classe Pydantic \texttt{Settings} (\texttt{app/core/config.py}) chargée depuis les variables d'environnement : aucun autre module ne lit \texttt{os.environ} directement, et les seuils peuvent être ajustés par déploiement sans modifier le code.

Le système emploie plusieurs patterns établis : \textbf{Singleton} pour les instances de modèles ML (chargées une fois par processus, injectées par dépendances FastAPI), \textbf{Repository} pour l'accès base de données via SQLAlchemy, \textbf{injection de dépendances} pour les handlers FastAPI (moquables en test), \textbf{Pipeline} pour le flux de modération (entrée $\to$ prédiction $\to$ routage $\to$ persistance $\to$ réponse) et \textbf{Stratégie} pour les fournisseurs de modèles image (CLIP, NudeNet, ensemble) sélectionnés par configuration.
